% GENERATED FILE. Edit canonical lessons, then run npm run generate:books.
% canonical-source-hash: fnv1a64:36861a4ff2935304
% canonical-lessons: AR-C123-battaniyya, AR-C123-wisada, AR-C123-sullam, AR-C123-mawqid, AR-C123-ina

\chapter{Blanket, Pillow, Ladder, Stove, Vessel}
\label{ch:ar-blanket-pillow-ladder-stove-vessel}
\hlchaptermodality{123}

\begin{chapteropening}
\textbf{By the end of this chapter:} \emph{I can say a blanket, a pillow, a ladder, a stove and a vessel.}

Complete the last lesson of chapter 123: I can say a blanket, a pillow, a ladder, a stove and a vessel.
\end{chapteropening}

\section[\texorpdfstring{\ar{بطانية} (baṭṭāniyya)}{بطانية (baṭṭāniyya)}]{\ar{بطانية} (baṭṭāniyya) — a blanket}

\label{lesson:AR-C123-battaniyya}



\begin{quote}
Before the new one: say the Arabic for a towel, then the Arabic for a mattress.
\end{quote}

\subsection*{You'll want to know: \ar{بطانية}}
\textbf{\ar{بطانية}} — \emph{baṭṭāniyya} — \textquotedblleft{}a blanket\textquotedblright{}.

\begin{cousinweb}
\textbf{\ar{بطانية}} (\emph{baṭṭāniyya}) is built on the native Arabic root \emph{b-ṭ-n}, \textquotedblleft{}the inside; a lining\textquotedblright{} (Lane). \emph{Baṭṭāniyya} is a modern word built on \emph{biṭāna}, \textquotedblleft{}a lining\textquotedblright{}.
\end{cousinweb}

\begin{grammarlens}[title={What you've built}]
One more word for this chapter.
\end{grammarlens}

\subsection*{Your turn}
\begin{itemize}
\raggedright
  \item \emph{Say it:} \emph{baṭṭāniyya}
  \item \emph{Say it:} \emph{baṭṭāniyya}, once more
  \item \emph{Say it:} say \emph{baṭṭāniyya}
  \item \emph{Recall:} say the Arabic for a towel, then the Arabic for a mattress, then say \emph{baṭṭāniyya} again
\end{itemize}

\begin{tcolorbox}[breakable,colback=teal!4,colframe=teal!35!black,title={Before you move on}]
What does \ar{بطانية} mean? (\textquotedblleft{}A blanket\textquotedblright{}.) Say it once more.
\end{tcolorbox}
\section[\texorpdfstring{\ar{وسادة} (wisāda)}{وسادة (wisāda)}]{\ar{وسادة} (wisāda) — a pillow}

\label{lesson:AR-C123-wisada}



\begin{quote}
Before the new one: say the Arabic for a mattress, then the Arabic for a blanket.
\end{quote}

\subsection*{You'll want to know: \ar{وسادة}}
\textbf{\ar{وسادة}} — \emph{wisāda} — \textquotedblleft{}a pillow\textquotedblright{}.

\begin{cousinweb}
\textbf{\ar{وسادة}} (\emph{wisāda}) is built on the native Arabic root \emph{w-s-d}, \textquotedblleft{}a pillow\textquotedblright{} (Lane). \emph{Wisāda} is a pillow, a cushion.
\end{cousinweb}

\begin{grammarlens}[title={What you've built}]
One more word for this chapter.
\end{grammarlens}

\subsection*{Your turn}
\begin{itemize}
\raggedright
  \item \emph{Say it:} \emph{wisāda}
  \item \emph{Say it:} \emph{wisāda}, once more
  \item \emph{Say it:} say \emph{wisāda}
  \item \emph{Recall:} say the Arabic for a mattress, then the Arabic for a blanket, then say \emph{wisāda} again
\end{itemize}

\begin{tcolorbox}[breakable,colback=teal!4,colframe=teal!35!black,title={Before you move on}]
What does \ar{وسادة} mean? (\textquotedblleft{}A pillow\textquotedblright{}.) Say it once more.
\end{tcolorbox}
\section[\texorpdfstring{\ar{سلم} (sullam)}{سلم (sullam)}]{\ar{سلم} (sullam) — a ladder}

\label{lesson:AR-C123-sullam}



\begin{quote}
Before the new one: say the Arabic for a blanket, then the Arabic for a pillow.
\end{quote}

\subsection*{You'll want to know: \ar{سلم}}
\textbf{\ar{سلم}} — \emph{sullam} — \textquotedblleft{}a ladder\textquotedblright{}.

\begin{cousinweb}
\textbf{\ar{سلم}} (\emph{sullam}) is built on the native Arabic root \emph{s-l-m}, \textquotedblleft{}a ladder\textquotedblright{} (Lane). Lane lists \emph{sullam} under \emph{s-l-m}; Hebrew \emph{sullam}, \textquotedblleft{}ladder\textquotedblright{}, is a cousin.
\end{cousinweb}

\begin{grammarlens}[title={What you've built}]
One more word for this chapter.
\end{grammarlens}

\subsection*{Your turn}
\begin{itemize}
\raggedright
  \item \emph{Say it:} \emph{sullam}
  \item \emph{Say it:} \emph{sullam}, once more
  \item \emph{Say it:} say \emph{sullam}
  \item \emph{Recall:} say the Arabic for a blanket, then the Arabic for a pillow, then say \emph{sullam} again
\end{itemize}

\begin{tcolorbox}[breakable,colback=teal!4,colframe=teal!35!black,title={Before you move on}]
What does \ar{سلم} mean? (\textquotedblleft{}A ladder\textquotedblright{}.) Say it once more.
\end{tcolorbox}
\section[\texorpdfstring{\ar{موقد} (mawqid)}{موقد (mawqid)}]{\ar{موقد} (mawqid) — a stove}

\label{lesson:AR-C123-mawqid}



\begin{quote}
Before the new one: say the Arabic for a pillow, then the Arabic for a ladder.

Say \textquotedblleft{}day.\textquotedblright{} (\emph{Yawm} — the \emph{yom} of \emph{Yom Kippur}.)
\end{quote}

\subsection*{You'll want to know: \ar{موقد}}
\textbf{\ar{موقد}} — \emph{mawqid} — \textquotedblleft{}a stove\textquotedblright{}.

\begin{cousinweb}
\textbf{\ar{موقد}} (\emph{mawqid}) is built on the native Arabic root \emph{w-q-d}, \textquotedblleft{}to burn, to kindle\textquotedblright{} (Lane). \emph{Mawqid} is the \emph{ma-} place noun, \textquotedblleft{}a place of kindling\textquotedblright{}.
\end{cousinweb}

\begin{grammarlens}[title={What you've built}]
One more word for this chapter.
\end{grammarlens}

\subsection*{Your turn}
\begin{itemize}
\raggedright
  \item \emph{Say it:} \emph{mawqid}
  \item \emph{Say it:} \emph{mawqid}, once more
  \item \emph{Say it:} say \emph{mawqid}
  \item \emph{Recall:} say the Arabic for a pillow, then the Arabic for a ladder, then say \emph{mawqid} again
\end{itemize}

\begin{tcolorbox}[breakable,colback=teal!4,colframe=teal!35!black,title={Before you move on}]
What does \ar{موقد} mean? (\textquotedblleft{}A stove\textquotedblright{}.) Say it once more.
\end{tcolorbox}
\section[\texorpdfstring{\ar{إناء} (ināʾ)}{إناء (ināʾ)}]{\ar{إناء} (ināʾ) — a vessel}

\label{lesson:AR-C123-ina}



\begin{quote}
Before the new one: say the Arabic for a ladder, then the Arabic for a stove.
\end{quote}

\subsection*{You'll want to know: \ar{إناء}}
\textbf{\ar{إناء}} — \emph{ināʾ} — \textquotedblleft{}a vessel\textquotedblright{}.

\begin{cousinweb}
\textbf{\ar{إناء}} (\emph{ināʾ}) is built on the native Arabic root \emph{ʾ-n-y}, \textquotedblleft{}a vessel\textquotedblright{} (Lane). Lane's \emph{ināʾ} is a vessel, a container.
\end{cousinweb}

\begin{grammarlens}[title={What you've built}]
One more word for this chapter.
\end{grammarlens}

\subsection*{Your turn}
\begin{itemize}
\raggedright
  \item \emph{Say it:} \emph{ināʾ}
  \item \emph{Say it:} \emph{ināʾ}, once more
  \item \emph{Say it:} say \emph{ināʾ}
  \item \emph{Recall:} say the Arabic for a ladder, then the Arabic for a stove, then say \emph{ināʾ} again
\end{itemize}

\begin{tcolorbox}[breakable,colback=teal!4,colframe=teal!35!black,title={Before you move on}]
What does \ar{إناء} mean? (\textquotedblleft{}A vessel\textquotedblright{}.) Say it once more.
\end{tcolorbox}
